\section{What ordinary training and other stability notions do not guarantee}
\label{sec:naturalness-obstructions}

The following results delimit the empirical question. They do not imply
that fragile outcomes are common under standard initialization, or that
their Boolean functions are statistically difficult.

\begin{proposition}[An open gradient-descent basin of fragile exact fits]
For every even $L\ge256$, there is a width-one, depth-$L$ tanh network
on $x\in\{-1,1\}$ with output $z(x)=x$, energy $\kappa_1\le L+76$,
and a postactivation sign-changing disturbance of size at most
\[
 \frac2L e^{-\sqrt L/4}.
\]
This disturbance is smaller than both calibrated support radii
\[
 r_{\rm ref}=\frac{.01}{\sqrt{1+(L-1)E_v}},\qquad
 r_S=\frac{.01}{\sqrt{1+S}}
\]
defined in Section~\ref{sec:model}.
For sufficiently small fixed step size, full-batch gradient descent on
the unregularized squared loss for the two labeled inputs has an
\emph{open set} of initializations converging to nearby exact fits that
still violate both support conditions.
\end{proposition}

\begin{proof}
Put $k_s=\min(s,L-s)$ and
\[
 r_s=\frac{\exp[-(k_s-1)/\sqrt L]}{2k_s},\quad
 a=\operatorname{atanh}(1/2),\quad
 b_s=\frac{\operatorname{atanh}(r_s)}{r_{s-1}},\quad v=2.
\]
Use first-layer input coefficient $a$ and bias zero, followed by scalar
weights $b_s$. Then $h_s(x)=r_sx$ and $z(x)=x$.
The reciprocal-ratio estimate from the bounded-excess construction in~\cite{previous} gives $\kappa_1\le L+76$; replacing its constant input by
$x$ preserves the computation and reduces the first-layer energy charge.
Also $b_s\ge e^{-1/2}/2>0.30$.

At $x=1$, the middle-layer disturbance $-2r_{L/2}$ changes the entire
remaining computation to its negative and gives final output $-1$.
Its norm is at most $2L^{-1}e^{-\sqrt L/4}$. Here $E_v=4$, so the ratio
of this upper bound to $r_{\rm ref}$ is at most
$400e^{-\sqrt L/4}/\sqrt L<1$ for $L\ge256$.
Since $S=\kappa_1\le L+76$ and $L+77\le4L-3$, it is also below $r_S$.

For the training claim, let $F(\theta)=(z_\theta(1),z_\theta(-1))$ and
$\ell(\theta)=\frac12\|F(\theta)-(1,-1)\|^2$.
At the constructed point $\theta_*$, the columns of $DF$ corresponding
to the first input coefficient and bias are $(c,-c)$ and $(c,c)$,
respectively, with $c>0$. Thus $DF$ has full row rank.
In a sufficiently small closed ball about $\theta_*$, its least singular
value is at least $\mu>0$, its norm is at most $G$, and
$\|\nabla^2\ell\|\le M$. The clean margins remain $>1/2$, and the
\emph{fixed} middle-layer attack remains strictly inside both calibrated
balls and gives strictly negative output.

For $0<\eta\le1/M$, the descent lemma gives, as long as the iterates
remain in this neighborhood,
\[
 \ell(\theta_{j+1})\le(1-\eta\mu^2)\ell(\theta_j).
\]
Their total movement is at most
\[
 \frac{\eta G\sqrt{2\ell(\theta_0)}}{1-\sqrt{1-\eta\mu^2}}
 \le \frac{2G\sqrt{2\ell(\theta_0)}}{\mu^2}.
\]
Starting in a sufficiently small open ball makes this displacement smaller
than the distance to the boundary, proving by induction that the iterates
stay inside. The summable increments give convergence, and the geometric
loss decay gives exact interpolation at the limit. Every limit retains the
strict fragility witness.
\end{proof}

A nondegenerate Gaussian initialization gives this open basin positive
probability. There is no useful lower bound here on that probability;
the result rules out a universal guarantee, not a typical-case guarantee.
It concerns unregularized full-batch squared-loss GD, rather than arbitrary
weight decay, Adam, or a prescribed finite training budget.

The same construction also violates the numerical-stability rule with
the radius in~\eqref{eq:numeric-radius} and permitted output change $1/16$.
For even $L\ge1024$, the attack-to-radius ratio is at most
$4000e^{-\sqrt L/4}/\sqrt L<1$, and the attack changes the output by two.
The same open-basin argument applies. Thus that alternative condition is
also not a universal consequence of ordinary gradient descent.

\begin{proposition}[Parameter flatness need not give activation robustness]
The networks just constructed preserve their correct signs under every
simultaneous Euclidean parameter perturbation of norm at most $1/4$.
At the exact fit their squared-loss Hessian has operator norm $O(L)$.
Nevertheless their absolute internal-noise radius is at most
$2L^{-1}e^{-\sqrt L/4}$.
\end{proposition}
\begin{proof}
Every internal scalar weight and the readout stay positive under the stated
parameter perturbation. At both inputs the first signed preactivation is
at least $\operatorname{atanh}(1/2)-\sqrt2/4>0$. Odd monotonic tanh
layers therefore retain the signs throughout the parameter ball.

For $t>0$, $t\tanh'(t)/\tanh(t)\le1$, since
$\sinh(t)\cosh(t)\ge t$. Each internal log-weight elasticity of the
positive output is consequently at most one. All scalar weights are
bounded below by a positive numerical constant, so every parameter
derivative at either input is bounded by a numerical constant.
There are $L+1$ parameters, and at interpolation
$\nabla^2\ell=DF^\top DF$; its norm is therefore $O(L)$.
\end{proof}

This is a separation between stability notions. Parameter flatness could
still support a different learning theorem, but it cannot be substituted
for absolute activation robustness in the existing proof.

\begin{proposition}[Eventual membership after readout amplification]
Every fixed network with nonzero outputs on the finite Boolean cube belongs
to either calibrated support after sufficiently large positive readout
amplification.
\end{proposition}
\begin{proof}
Continuity and finiteness of the input set give $\rho,m>0$ such that every
hidden disturbance of total normalized norm at most $\rho$ preserves each
clean sign and leaves unscaled signed output at least $m$.
Multiplying the readout by $A$ makes either calibrated radius $r_A$ tend
to zero. For large enough $A$, $r_A<\rho$ and $Am>r_A$, so even adding
the permitted scalar output noise cannot change a sign. Increasing $A$
also makes the clean output margin exceed $1/2$.
\end{proof}

This does not remove the Gaussian cost: the readout energy grows as $A^2$,
and the required $A$ may be exponentially large. It does mean that empirical
membership before and after output rescaling answers different questions.
The calibrated condition is a joint restriction on robustness and scale.

\begin{proposition}[Global and typical margins can differ exponentially]
There is a one-hidden-unit network with bounded input weight norm and
Boolean function $x_1$ whose minimum output margin is $\tanh(2^{-d})$.
Enforcing the original global clean margin $>1/2$ by readout amplification
requires a multiplier at least $2^{d-1}$, while permitting an exceptional
fraction of inputs requires only polynomial amplification in the reciprocal
exceptional fraction.
\end{proposition}
\begin{proof}
Take
\[
 a_d(x)=\sum_{j=1}^d2^{-j}x_j,\qquad z_d(x)=\tanh(a_d(x)).
\]
The squared input weight norm is less than $1/3$. The first coefficient
exceeds the sum of the others, so the sign is $x_1$, and the minimum
absolute preactivation is $2^{-d}$. Thus the global clean-margin multiplier
must exceed $[2\tanh(2^{-d})]^{-1}\ge2^{d-1}$.
Under the uniform cube, $a_d$ is uniform on the odd multiples of $2^{-d}$
in $(-1,1)$, giving
\[
 \Pr(|a_d(X)|\le u)\le u+2^{-d}\qquad(0\le u\le1).
\]
Consequently for $A\ge1$,
\[
 \Pr(|Az_d(X)|\le1/2)
 \le\operatorname{atanh}(1/(2A))+2^{-d}
 \le A^{-1}+2^{-d}.
\]
For example, $A\ge2/\epsilon$ suffices when $2^{-d}\le\epsilon/2$.
\end{proof}

The last example concerns the inherited global clean-margin convention,
not an additional failure of internal robustness. It motivates distinguishing
training, held-out, and worst-input quantities, and investigating support
conditions that explicitly tolerate a small exceptional input mass.

\begin{proposition}[A generic Gaussian halfspace has a tiny cube margin]
Let $d\ge64$ and $w\sim\mathcal N(0,I_d/d)$. Then
\[
 \Pr\left\{\min_{x\in\{-1,1\}^d}|w\cdot x|
                    \le d^2 2^{-d}\right\}\ge1-\frac6d.
\]
On this event, the one-hidden-unit output $A\tanh(w\cdot x)$ needs
$A>2^{d-1}/d^2$ to satisfy the global clean margin $>1/2$.
This is a statement about Gaussian coefficients, not the distribution of
weights after training and not a sample-complexity lower bound.
\end{proposition}
\begin{proof}
For $0<\tau\le d^{-1/2}$ put
$Z=2^{-d}\sum_x\mathbf1\{|w\cdot x|\le\tau\}$ and
$p=\mathbb EZ=\Pr(|G|\le\tau)$ for $G\sim\mathcal N(0,1)$.
For independent uniform cube points $x,x'$, their Gaussian evaluations
have correlation $\rho=d^{-1}\sum_jx_jx'_j$.
When $\rho\ne\pm1$, the bivariate Gaussian density bound gives
\[
 \frac{\Pr(|w\cdot x|\le\tau,\ |w\cdot x'|\le\tau)}{p^2}
 \le\frac{e^{\tau^2}}{\sqrt{1-\rho^2}},
\]
using $p\ge\sqrt{2/\pi}\,\tau e^{-\tau^2/2}$.
The antipodal or identical pairs have total probability $2^{1-d}$.
For $|\rho|\le1/2$, $(1-\rho^2)^{-1/2}\le1+\rho^2$.
Otherwise, excluding the two degenerate correlations, this factor is at
most $\sqrt{d/2}$. Since $\mathbb E\rho^2=1/d$ and
$\Pr(|\rho|>1/2)\le2e^{-d/8}$,
\[
 \mathbb E\!\left[
   \frac{\mathbf1\{\rho\ne\pm1\}}{\sqrt{1-\rho^2}}\right]
 \le1+\frac1d+\sqrt{2d}\,e^{-d/8}\le1+\frac2d.
\]
Here the last inequality holds for $d\ge64$.
Using $p\ge\tau/2$ and $e^{\tau^2}\le1+2/d$, Chebyshev gives
\[
 \Pr(Z=0)\le\frac{\operatorname{Var}Z}{p^2}
 \le\frac{2^{2-d}}\tau+\frac5d.
\]
Take $\tau=d^22^{-d}$; this is at most $d^{-1/2}$ for $d\ge64$.
The right-hand side is at most $6/d$. The output-scaling claim follows
from $|\tanh u|\le|u|$.
\end{proof}

The same observation includes a Gaussian first-layer bias by passing to
the $(d+1)$-dimensional cube: absolute affine values on the original cube
are exactly absolute homogeneous values on one representative from every
antipodal pair of the augmented cube.
